\documentclass[11pt]{article}
\usepackage{ochamath}
\subjectcolor{2F5DA8}
\setsheet{線形代数・電気系院試補充}{特性多項式と固有値の重複　演習}{https://university-math-crj.pages.dev/linear-algebra/characteristic-polynomial/}
\namelinetrue
\begin{document}
\makesheethead{問題}
自作問題。本文の例題とは数値や設定が異なります。条件・途中式・検算を示してください。
\problem[重根と固有空間]
$A=\begin{pmatrix}4&2&0\\0&4&0\\0&0&-1\end{pmatrix}$ の特性多項式、各重複度、対角化可能性を求めよ。
\workspace{45mm}
\problem[パラメータと対角化]
$A(t)=\begin{pmatrix}3&t&0\\0&3&0\\0&0&1\end{pmatrix}$ が実対角化できる $t$ を求めよ。
\workspace{45mm}
\newpage
\problem[3次の係数]
固有値が $1,2,4$ の3次行列について特性多項式、跡、行列式、$\operatorname{tr}(A^2)$ を求めよ。
\workspace{45mm}
\problem[同じ多項式でも相似でない]
$D=\operatorname{diag}(5,5)$ と $J=\begin{pmatrix}5&1\\0&5\end{pmatrix}$ が相似でないことを証明せよ。
\workspace{45mm}
\newpage
\problem[体による違い]
$R=\begin{pmatrix}0&-2\\2&0\end{pmatrix}$ の固有値を求め、実数と複素数で対角化可能性を答えよ。
\workspace{45mm}
\problem[固有空間の独立性]
相異なる固有値の固有空間の和が直和であることを証明せよ。
\workspace{45mm}
\end{document}
